\section*{الفصل \ref{ch:b1:extent-q-and-k} --- وصف نظام كيميائي: التقدم والنشاطيات وخارج التفاعل وثابت التوازن}

\begin{solution}{exo:b1:extent-q-and-k:1}
\omterm{def:b1:extent-q-and-k:variables}{المكثفة}: درجة الحرارة، \omterm{def:b1:extent-q-and-k:composition}{والتركيز المولي}، والكتلة الحجمية، والضغط، \omterm{def:b1:extent-q-and-k:composition}{والكسر
المولي}. \omterm{def:b1:extent-q-and-k:variables}{والممتدة}: الحجم، وكمية المادة، والكتلة.
\end{solution}

\begin{solution}{exo:b1:extent-q-and-k:2}
$p(\ce{N2}) = 0.7808 \times 1.013 = \qty{0.791}{bar}$، $p(\ce{O2}) =
0.2095 \times 1.013 = \qty{0.212}{bar}$، $p(\ce{Ar}) = 0.0093 \times 1.013
= \qty{0.0094}{bar}$.
\end{solution}

\begin{solution}{exo:b1:extent-q-and-k:3}
$n(\ce{H2}) = 3.0 - 2\xi$، $n(\ce{O2}) = 2.0 - \xi$، $n(\ce{H2O}) =
2\xi$. $\xi_{\max} = \min(3.0/2,\ 2.0/1) = \qty{1.5}{mol}$: فثنائي الهيدروجين هو
المحدِّد. الحالة النهائية: \ce{H2} 0، \ce{O2} \qty{0.5}{mol}، \ce{H2O}
\qty{3.0}{mol}.
\end{solution}

\begin{solution}{exo:b1:extent-q-and-k:4}
(أ) $Q = p_{\ce{CO2}}/p^\circ$. (ب) $Q = (c^\circ)^2/([\ce{Ag+}][\ce{Cl-}])$.
(ج) $Q = (p_{\ce{NH3}}/p^\circ)^2/[(p_{\ce{N2}}/p^\circ)(p_{\ce{H2}}/p^\circ)^3]$.
(د) $Q = [\ce{NH4+}][\ce{OH-}]/([\ce{NH3}]\,c^\circ)$، إذ \omterm{def:b1:extent-q-and-k:activity}{نشاطية} الماء 1.
\end{solution}

\begin{solution}{exo:b1:extent-q-and-k:5}
$\mathrm{A \rightleftharpoons C}$ هو المجموع: $K = 10 \times 0.5 = 5$. $\mathrm{C \rightleftharpoons A}$: $1/5 =
0.2$. $\mathrm{2\,A \rightleftharpoons 2\,C}$: $5^2 = 25$.
\end{solution}

\begin{solution}{exo:b1:extent-q-and-k:6}
$Q = 5.0 \times 10^{-5} < K^\circ$: الاتجاه المباشر. $Q = 0.10 > K^\circ$: الاتجاه المعاكس.
متفاعلات فقط: $Q = 0 < K^\circ$، الاتجاه المباشر.
\end{solution}

\begin{solution}{exo:b1:extent-q-and-k:7}
$n(\ce{A2}) = 1 - \xi$، $n(\ce{A}) = 2\xi$، والمجموع $1 + \xi$. ومع $p =
p^\circ$، $Q = \dfrac{(2\xi/(1+\xi))^2}{(1-\xi)/(1+\xi)} = \dfrac{4\xi^2}{1 -
\xi^2}$. و$Q = 0.25$ يعطي $4.25\,\xi^2 = 0.25$، $\xi = \qty{0.243}{mol}$.
\end{solution}

\begin{solution}{exo:b1:extent-q-and-k:8}
مع $x = [\ce{C}]$: $x/(0.10 - x)^2 = 100$، فيكون $100x^2 - 21x + 1 = 0$
و$x = (21 - \sqrt{41})/200 = \qty{0.0730}{mol/L}$ (الجذر الآخر يتجاوز
0.10). $[\ce{A}] = [\ce{B}] = \qty{0.0270}{mol/L}$، $[\ce{C}] =
\qty{0.0730}{mol/L}$، $\tau = 0.73$.
\end{solution}

\begin{solution}{exo:b1:extent-q-and-k:9}
انطلاقًا من $n$ من كل منهما: $Q = \xi^2/(n - \xi)^2 = K^\circ$، فيكون $\xi/(n - \xi) =
\sqrt{K^\circ}$ و$\tau = \xi/n = \sqrt{K^\circ}/(1 + \sqrt{K^\circ})$.
ويقتضي $\tau = 0.99$ أن يكون $\sqrt{K^\circ} = 99$، أي $K^\circ \approx \num{9.8e3}$:
نحو $10^4$، وهي العتبة المعتادة للتفاعل الكمّي.
\end{solution}

\begin{solution}{exo:b1:extent-q-and-k:10}
عند التوازن $p_{\ce{CO2}} = 0.20\,p^\circ$، أي $n(\ce{CO2}) =
pV/RT = 0.20 \times 10^5 \times 0.0100/(8.314 \times 1100) =
\qty{0.0219}{mol}$. ومع \qty{0.010}{mol} من الكربونات لا يمكن بلوغ ذلك:
فتتفكك كلها، $p_{\ce{CO2}} = 0.010 \times 8.314 \times
1100/0.0100 = \qty{9.15e3}{Pa} = \qty{0.091}{bar}$، و$Q < K^\circ$،
والحالة النهائية (\qty{0.010}{mol} من \ce{CaO}، و\qty{0.010}{mol} من \ce{CO2}، ولا
\ce{CaCO3}) ليست \omterm{def:b1:extent-q-and-k:quantitative}{حالة توازن}. ومع \qty{0.050}{mol}: توازن،
\qty{0.0219}{mol} من \ce{CO2} ومن \ce{CaO}، و\qty{0.0281}{mol} من \ce{CaCO3} باقية.
\end{solution}

\begin{solution}{exo:b1:extent-q-and-k:11}
عند التوازن $n(\ce{B}) = 2\,n(\ce{A})$ و$n(\ce{C}) = 3\,n(\ce{A})$
(يُختصر الحجم). ومنه $6\,n(\ce{A}) = 1.0$: $n(\ce{A}) =
\qty{0.167}{mol}$، $n(\ce{B}) = \qty{0.333}{mol}$، $n(\ce{C}) =
\qty{0.500}{mol}$.
\end{solution}

\begin{solution}{exo:b1:extent-q-and-k:12}
كميتان متساويتان: $\xi^2/(0.50 - \xi)^2 = 4$، $\xi/(0.50 - \xi) = 2$،
$\xi = \qty{0.333}{mol}$، \omterm{def:b1:extent-q-and-k:fractional-extent}{والمردود} $0.333/0.50 = 67\,\%$. ومع \qty{1.0}{mol}
من الكحول: $\xi^2 = 4(0.50 - \xi)(1.0 - \xi)$، $3\xi^2 - 6\xi + 2 = 0$،
$\xi = (6 - \sqrt{12})/6 = \qty{0.423}{mol}$، \omterm{def:b1:extent-q-and-k:fractional-extent}{والمردود} 85\,\%: فإفراط
أحد المتفاعلين يرفع \omterm{def:b1:extent-q-and-k:fractional-extent}{مردود} ناتج الآخر.
\end{solution}

\begin{solution}{pb:b1:extent-q-and-k:1}
\textbf{1.} $x(\ce{CH4}) = 1.00/4.00 = 0.250$؛ $x(\ce{H2O}) = 0.750$.
\textbf{2.} \qty{7.5}{bar} و\qty{22.5}{bar}.
\textbf{3.} الكتلتان 16.04 و\qty{54.06}{g}: $w(\ce{CH4}) = 0.229$،
$w(\ce{H2O}) = 0.771$.
\textbf{4.} \omterm{def:b1:extent-q-and-k:composition}{الكسور المولية}، \omterm{def:b1:extent-q-and-k:composition}{والكسور الكتلية}، \omterm{def:b1:extent-q-and-k:composition}{والضغوط الجزئية} (و
الضغط الكلي)؛ أما الكميات \omterm{def:b1:extent-q-and-k:variables}{فممتدة}.
\textbf{5.} $a = p_i/p^\circ$: 7.5 للميثان، و22.5 لبخار الماء.
\textbf{6.} \ce{CH4} $1.00 - \xi$، \ce{H2O} $3.00 - \xi$، \ce{CO} $\xi$،
\ce{H2} $3\xi$، والمجموع $4.00 + 2\xi$.
\textbf{7.} $\xi_{\max} = \qty{1.00}{mol}$؛ والميثان هو المحدِّد.
\textbf{8.} \ce{CH4} 0، \ce{H2O} \qty{2.00}{mol}، \ce{CO} \qty{1.00}{mol}،
\ce{H2} \qty{3.00}{mol}.
\textbf{9.} \qty{6.00}{mol}، مقابل \qty{4.00}{mol} في التغذية: فالتفاعل
يصنع أربعة جزيئات من جزيئين.
\textbf{10.} \ce{H2O} 0.333، \ce{CO} 0.167، \ce{H2} 0.500.
\textbf{11.} لضمان تفاعل الميثان كله ولإبقاء بخار ماء
لتفاعل الإزاحة، الذي يخفض خارجَه.
\textbf{12.} \ce{CO} $1 - \xi$، \ce{H2O} $2 - \xi$، \ce{CO2} $\xi$، \ce{H2}
$3 + \xi$، والمجموع 6.00 عند كل تقدم.
\textbf{13.} $Q = \dfrac{(p_{\ce{CO2}}/p^\circ)(p_{\ce{H2}}/p^\circ)}
{(p_{\ce{CO}}/p^\circ)(p_{\ce{H2O}}/p^\circ)}$ مع $p_i = (n_i/6)p$: فتُختصر
العوامل $p/(6p^\circ)$ بين البسط والمقام (جزيئان غازيان
في كل طرف)، فتبقى العبارة المعطاة.
\textbf{14.} $\xi = 0$: $Q = 0 < 4.0$، الاتجاه المباشر.
\textbf{15.} $\xi(3 + \xi) = 4(1 - \xi)(2 - \xi)$، أي $3\xi^2 - 15\xi +
8 = 0$: $\xi = (15 - \sqrt{129})/6 = \qty{0.607}{mol}$؛ والجذر الآخر،
4.39، يتجاوز $\xi_{\max} = 1$.
\textbf{16.} \ce{CO} \qty{0.393}{mol}، \ce{H2O} \qty{1.393}{mol}،
\ce{CO2} \qty{0.607}{mol}، \ce{H2} \qty{3.607}{mol}.
\textbf{17.} $\tau = 0.607/1.00 = 0.607$.
\textbf{18.} $\xi(3 + \xi) = 4(1 - \xi)(5 - \xi)$، $3\xi^2 - 27\xi + 20 =
0$، $\xi = \qty{0.814}{mol}$: بخار أكثر، وتحويل أكثر.
\textbf{19.} $K^\circ = (0.99 \times 3.99)/(0.01 \times 1.01) = 391$.
\textbf{20.} الدالة $Q(\xi)$ متزايدة، فيُبلغ ثابت $K^\circ$ أكبر عند
تقدم أكبر. وثابت هذا التفاعل أكبر عند درجة
حرارة أدنى، حيث يكون التفاعل بطيئًا: فيعمل مفاعل أول ساخنًا
وسريعًا، ويُتم مفاعل ثانٍ أبرد التحويل (أما تعلق
$K^\circ$ بدرجة الحرارة فيُعالج في مجلد السنة الثانية).
\textbf{21.} $3.607/(6.00 - 1.393) = 3.607/4.607 = 0.783$.
\textbf{22.} \qty{0.393}{mol} من \ce{CO} و\qty{0.607}{mol} من \ce{CO2}.
\textbf{23.} $x(\ce{H2}) = 3.607/6.00 = \textbf{0.60}$: ستة جزيئات من كل
عشرة تغادر مفاعل الإزاحة هي جزيئات هيدروجين.
\end{solution}
